\chapter{Introduction}
The invasion of Ukraine perpetrated by the Russian Federation is one of the most important political events of the 3rd decade of the 21\textsuperscript{st} century. Since the beginning
of the invasion approximately 14600 civilians have died and the total number of people who lost their lives is approximately of 2 millions as of \citet{jonesrussia}. This war has been particularly historical
because has been
the first major conflict between two different european states since the end of the Second World War,
 as the war in Iugoslavia is considered as civil conflict.\\
The financial and military aid of the USA, the European Union and the United Kingdom to Ukraine played a major role in 
letting the state to resist to the invasion and, for this reason,
the support of these countries' public opinions in favour or against the war could lead to major results on the battlefield. 
Consequently, both Western countries and the 
Russian Federation have started a huge campaign of propaganda to influence the public opinions of the countries who are supporting Ukraine, therefore an automatic system
which allows to understand the the major topics which are discussed in the society regarding the war 
in Ukraine could reveal to be fundamental for winning on the frontline.
\section{ANR CIGAIA project}
As mentioned on its website, the CIGAIA project has a dual nature as it aims at understanding the argumentative patterns
used in the civil controversies about international matterss, but also allowing to better understand the information environment 
of military operations.\\
This thesis aims at working on what the project defines as the second axis of the project: designing an Artificial Intelligence system
for automatically analysing the the argumentation in these controversies, as a first passage which will allow in the future 
to design a decision support tool. The goal of the latter tool will be to map the actors via the identification of controversies
using argumentation and counter-argumentation strategies.
\section{Argument Mining}\label{sec:am}
Argument Mining (AM) is one of the most interesting topics currently in the field of Natural Language Processing and it was 
defined by \citet{palau}
as a study which \textquote{aims to detect the arguments presented in a text document, the relations between them and the internal structure of each individual argument}. In this context three main 
tasks were defined to address the various aspects which argumentation involves:
\begin{itemize}
    \item argument classification (AC), where the goal is to isolate the parts which are considered argumentative from the ones which are deemed non-argumentative;
    \item argument component classification (ACC), where the objective is to identify and classify argument components as either claims, assertions of a given fact which the speaker wants to prove in the text, or premises, 
    axioms which are assumed to be true in the text;
    \item argument relation classification (ARC), where the goal is to identify and classify different types of relationship between claims and premises (i.e. a premise can be used to support a claim but also to contest it).
\end{itemize}
In this project it will be assumed that the task of segmentation, therefore the one isolating the senteces which are deemend of interest for classification,
was solved. For this reason the focus of this project will be on the classification part of the 3 subtasks.\\
After this brief introduction, chapter \ref{chapter:2} will focus on the related work and describing the state-of-the-art for classification in argument mining, chapter \ref{chapter:3}
will expose the methodology which will be employed in this project, chapter \ref{chapter:4} the results of hte project and chapter \ref{chapter:5} will contain the conlcusions and the 
limitations obtained.